\documentclass[UTF8,fontset=fandol,a4paper,10pt]{ctexart}
\usepackage[margin=22mm]{geometry}
\usepackage[colorlinks=true,urlcolor=blue]{hyperref}
\setlength{\parindent}{0pt}
\setlength{\parskip}{6pt}
\title{徐乔依}
\author{Joyce}
\date{}
\begin{document}
\maketitle
\section*{关于我}
我叫 Qiaoyi Xu（Joyce），家乡是北京。目前在北京大学攻读流行病学与卫生统计学博士学位，此前在哥伦比亚大学和马萨诸塞大学阿默斯特分校学习公共卫生、流行病学与数学。我的职业目标是将所学用于健康领域，帮助人们改善健康。
\section*{教育背景}
\textbf{北京大学 · Beijing, China}

2026 年 9 月至预计 2030 年 7 月，Doctor of Philosophy in Epidemiology and Biostatistics。

\textbf{哥伦比亚大学 · New York, NY}

2021 年 9 月至 2023 年 5 月，Master of Public Health in Epidemiolog, Certificate in Applied Biostatistics。

GPA：3.93。

\textbf{马萨诸塞大学阿默斯特分校 · Amherst, MA}

2018 年 9 月至 2021 年 5 月，Bachelor of Science in Public Health Sciences and Mathematics, Double Major。

GPA：3.8。

\section*{工作与实习经历}
\begin{itemize}
\item \textbf{Pharmaron Clinical · Biostatistician}（2023 年 7 月至 2026 年 8 月）：参与 I/II 期临床试验、生物等效性研究、研究者发起试验及 FDA IND 申请支持分析；编写 SAP、TLF 模板等统计文件，审核研究方案与相关交付文件。
\item \textbf{Pfizer · 流行病学研究助理}（2022 年 6–8 月，纽约）：支持莱姆病疫苗相关流行病学研究，开展加拿大蜱监测系统文献综述，使用 SAS 分析数据并汇报结果。
\item \textbf{Frost \& Sullivan · 咨询分析实习生}（2021 年 12 月至 2022 年 2 月，上海）：参与资料收集、市场分析展示及专家访谈纪要整理。
\item \textbf{National Center for Medical Service Administration · 数据分析实习}（2020 年 7–9 月，北京）：收集核酸检测数据、开展分析并撰写日报。
\item \textbf{中国疾病预防控制中心 · 数据分析与健康传播实习}（2019 年 6–8 月，北京）：收集和分析环境健康相关数据，参与环境健康科普活动。
\end{itemize}
\section*{数据分析作品}
纽约州 NOAA 气象数据探索。
\href{https://joycexu2000-web.github.io/personal-homepage-python/plotly.html}{Plotly} · \href{https://joycexu2000-web.github.io/personal-homepage-python/dashboard.html}{Dashboard}
\section*{技能与兴趣}
分析工具：R、SAS、MySQL、Excel。\\ 语言：中文、英文。\\ 个人兴趣：舞蹈、冲浪、百老汇音乐剧；我的猫叫 Erer。
\section*{联系与资料}
Phone number: +86 15711312382\\ Email: \href{mailto:qiaoyixu2000@gmail.com}{qiaoyixu2000@gmail.com}\\ Address: Beijing, China
\end{document}